% =========================================================
% Isometry preserves diameter — explanation + diagram
% =========================================================
% Requires:
% \usepackage{amsmath,amssymb}
% \usepackage{tikz}
% \usetikzlibrary{arrows.meta,calc,decorations.pathmorphing,positioning}

\subsubsection*{Why it is OK that $A\subseteq X$ and $f(A)\subseteq Y$}
Even though $A$ and $f(A)$ live in different metric spaces, the \emph{diameter}
comparison does not require comparing points across spaces.  Instead, it compares
two sets of \emph{real numbers} (distance values).

Let $(X,d_X)$ and $(Y,d_Y)$ be metric spaces and let $f:X\to Y$ be an isometry, i.e.
\[
d_Y\bigl(f(x),f(y)\bigr)=d_X(x,y)\qquad\text{for all }x,y\in X.
\]
Fix $A\subseteq X$ and consider the corresponding \emph{distance sets}
\[
D_A := \{\, d_X(x,y) : x,y\in A \,\}\subseteq \mathbb{R},
\qquad
D_{f(A)} := \{\, d_Y(u,v) : u,v\in f(A) \,\}\subseteq \mathbb{R}.
\]
These two sets live in the \emph{same ambient space} $\mathbb{R}$, so it makes
perfect sense to compare them. Moreover, by the definition of image,
every $u\in f(A)$ has the form $u=f(x)$ for some $x\in A$, and similarly for $v$.
Thus each distance in $D_{f(A)}$ can be written as $d_Y(f(x),f(y))$ with $x,y\in A$,
and the isometry property gives
\[
d_Y\bigl(f(x),f(y)\bigr)=d_X(x,y).
\]
So the set of distances realized by pairs of points in $A$ is the same as the set of
distances realized by pairs of points in $f(A)$; consequently their suprema agree:
\[
\operatorname{diam}(A)=\sup D_A,
\qquad
\operatorname{diam}(f(A))=\sup D_{f(A)}.
\]
In short: \emph{points live in $X$ and $Y$, but distance values live in $\mathbb{R}$.}

% ---------------------------------------------------------
% TikZ diagram: two metric spaces, subset A and f(A),
% points x1,x2 and f(x1),f(x2), and distance values in R.
% ---------------------------------------------------------

\begin{center}
\begin{tikzpicture}[
  >=Latex,
  every node/.style={font=\small},
  blob/.style={draw, thick, rounded corners=20pt},
  setfillA/.style={fill=red!6, draw=red!70!black, thick},
  setfillB/.style={fill=blue!6, draw=blue!70!black, thick},
  map/.style={->, thick, blue!70!black},
  dist/.style={thick, teal!70!black},
  pt/.style={circle, fill=black, inner sep=1.1pt}
]

% --- ambient "spaces" as blobs
\draw[blob, red!55!black]  (-5.6,-1.9) rectangle (-1.4,1.9);
\node[red!55!black] at (-3.5,1.65) {$X$};

\draw[blob, blue!55!black] ( 1.4,-1.9) rectangle ( 5.6,1.9);
\node[blue!55!black] at (3.5,1.65) {$Y$};

% --- subset A inside X (a smaller blob)
\draw[setfillA]
  (-5.05,-0.95)
  .. controls (-5.25,0.10) and (-4.55,0.95) .. (-3.60,0.90)
  .. controls (-2.70,0.85) and (-2.20,0.20) .. (-2.25,-0.55)
  .. controls (-2.30,-1.25) and (-3.50,-1.35) .. (-4.30,-1.10)
  .. controls (-4.70,-0.98) and (-4.85,-1.05) .. (-5.05,-0.95)
  -- cycle;
\node[red!70!black] at (-3.75,1.05) {$A$};

% --- subset f(A) inside Y (a smaller blob)
\draw[setfillB]
  (2.00,-0.95)
  .. controls (1.80,0.10) and (2.50,0.95) .. (3.45,0.90)
  .. controls (4.35,0.85) and (4.85,0.20) .. (4.80,-0.55)
  .. controls (4.75,-1.25) and (3.55,-1.35) .. (2.75,-1.10)
  .. controls (2.35,-0.98) and (2.20,-1.05) .. (2.00,-0.95)
  -- cycle;
\node[blue!70!black] at (3.55,1.05) {$f(A)$};

% --- points in A
\node[pt] (x1) at (-4.55,0.25) {};
\node[pt] (x2) at (-3.05,-0.35) {};
\node[above left=1pt of x1] {$x_1$};
\node[below right=1pt of x2] {$x_2$};

% --- points in f(A)
\node[pt] (fx1) at (2.55,0.10) {};
\node[pt] (fx2) at (4.10,-0.25) {};
\node[above left=1pt of fx1] {$f(x_1)$};
\node[below right=1pt of fx2] {$f(x_2)$};

% --- distance segments
\draw[dist] (x1) -- (x2);
\node[teal!70!black] at ($(x1)!0.5!(x2)+(0.15,0.25)$)
  {$d_X(x_1,x_2)$};

\draw[dist] (fx1) -- (fx2);
\node[teal!70!black] at ($(fx1)!0.5!(fx2)+(0.00,0.25)$)
  {$d_Y(f(x_1),f(x_2))$};

% --- map arrows from points
\draw[map, bend left=15] (x1) to node[above, font=\scriptsize] {$f$} (fx1);
\draw[map, bend left=10] (x2) to node[below, font=\scriptsize] {$f$} (fx2);

% --- equality note (distance values live in R)
\node[align=center, font=\small] (Rnote) at (0,-1.45)
  {$d_Y(f(x_1),f(x_2))=d_X(x_1,x_2)\in\mathbb{R}$};

% --- diameter note
\node[align=center, font=\scriptsize] at (0,-2.25)
  {$\operatorname{diam}(A)=\sup\{d_X(x,y):x,y\in A\}$,\quad
   $\operatorname{diam}(f(A))=\sup\{d_Y(u,v):u,v\in f(A)\}$};

\end{tikzpicture}
\end{center}









% =========================================================
% Isometries Preserve Metric Structure
% =========================================================

\begin{theorem}[Isometries Preserve Metric Structure]
Let $(X,d_X)$ and $(Y,d_Y)$ be metric spaces and let
$f:X\to Y$ be an isometry, meaning
\[
d_Y(f(x),f(y)) = d_X(x,y)
\quad
\text{for all } x,y\in X.
\]

Then every property defined purely in terms of the metric
is preserved under $f$. In particular:

\begin{enumerate}

\item Distances are preserved:
\[
d_Y(f(x),f(y)) = d_X(x,y).
\]

\item Diameter is preserved:
\[
\operatorname{diam}(f(A)) = \operatorname{diam}(A)
\quad
\text{for every } A\subseteq X.
\]

\item Bounded sets remain bounded.

\item If $(x_n)$ is a Cauchy sequence in $X$, then
$(f(x_n))$ is Cauchy in $Y$.

\item If $x_n \to x$ in $X$, then
\[
f(x_n) \to f(x)
\quad
\text{in } Y.
\]

\end{enumerate}

\end{theorem}


% =========================================================
% Proof: Isometries Preserve Metric Structure
% =========================================================

\subsection*{Isometries Preserve Metric Structure}

\begin{proof}

\Claim Let $(X,d_X)$ and $(Y,d_Y)$ be metric spaces and let
$f:X\to Y$ be an isometry. Then every property defined purely
in terms of the metric is preserved under $f$. In particular,
for every $A\subseteq X$:

\begin{enumerate}
\item Distances are preserved.
\item Diameter is preserved.
\item Boundedness is preserved.
\item Cauchy sequences are preserved.
\item Convergence of sequences is preserved.
\end{enumerate}

\Given Metric spaces $(X,d_X)$ and $(Y,d_Y)$ and an isometry
$f:X\to Y$ satisfying
\[
d_Y(f(x),f(y)) = d_X(x,y)
\quad
\text{for all } x,y\in X.
\]

\Goal To show that metric properties defined using distances
are preserved under $f$.

\Strategy Each property listed above is defined using only the
metric. Since $f$ preserves all pairwise distances, the defining
conditions remain unchanged when points are mapped by $f$.

\medskip

\textbf{(1) Distances.}

By definition of isometry,
\[
d_Y(f(x),f(y)) = d_X(x,y)
\quad
\text{for all } x,y\in X.
\]

\medskip

\textbf{(2) Diameter.}

Recall
\[
\operatorname{diam}(A)
=
\sup\{d_X(x,y):x,y\in A\}.
\]

Since
\[
d_Y(f(x),f(y)) = d_X(x,y),
\]

we obtain
\[
\{d_Y(f(x),f(y)):x,y\in A\}
=
\{d_X(x,y):x,y\in A\}.
\]

Taking suprema yields
\[
\operatorname{diam}(f(A))=\operatorname{diam}(A).
\]

\medskip

\textbf{(3) Boundedness.}

A set $A$ is bounded if
\[
\sup\{d_X(x,y):x,y\in A\}<\infty.
\]

Since all distances are preserved by $f$, the same supremum
bounds distances in $f(A)$, so $f(A)$ is bounded whenever $A$
is bounded.

\medskip

\textbf{(4) Cauchy sequences.}

Suppose $(x_n)$ is Cauchy in $X$. Then

\[
\forall \varepsilon>0\;\exists N\;
\text{such that }
d_X(x_m,x_n)<\varepsilon
\quad
\text{for all } m,n\ge N.
\]

Using the isometry property,

\[
d_Y(f(x_m),f(x_n))
=
d_X(x_m,x_n)
<
\varepsilon
\]

for all $m,n\ge N$, showing that $(f(x_n))$ is Cauchy in $Y$.

\medskip

\textbf{(5) Convergence.}

Suppose $x_n\to x$ in $X$. Then

\[
d_X(x_n,x)\to 0.
\]

Using the isometry property,

\[
d_Y(f(x_n),f(x))
=
d_X(x_n,x).
\]

Hence

\[
d_Y(f(x_n),f(x))\to 0,
\]

which means

\[
f(x_n)\to f(x)
\quad
\text{in } Y.
\]

\end{proof}

